% !TEX root = ../algebra_02.tex

\section{Критерий диагонализируемости в терминах геометрических и алгебраических кратностей}

\begin{Theorem}{Критерий через алгебраические кратности}{}
	Пусть $\alpha\in\operatorname{End} V$, $\dim V=n$. Оператор $\alpha$ диагонализируем тогда и только тогда, когда выполнены два условия:
	\begin{enumerate}
		\item характеристический многочлен $\chi_\alpha$ раскладывается над полем $K$ на линейные множители;
		\item для каждого собственного значения $\lambda$ выполнено равенство кратностей $g_\lambda=a_\lambda$.
	\end{enumerate}
\end{Theorem}

\begin{proof}
	Пусть $\alpha$ диагонализируем. Тогда в некотором базисе его матрица диагональна, и на диагонали стоят собственные значения $\lambda_i$ с кратностями $m_i$. Тогда характеристический многочлен имеет вид $\chi_\alpha(x) = (\lambda_1-x)^{m_1}\cdots(\lambda_k-x)^{m_k}$. Значит, он раскладывается на линейные множители. Кроме того, $g_{\lambda_i}=m_i$ (размерности соответствующих собственных подпространств), и $a_{\lambda_i}=m_i$ по виду характеристического многочлена. Следовательно, $g_{\lambda_i}=a_{\lambda_i}$ для всех $i$.

	Обратно, пусть $\chi_\alpha$ раскладывается на линейные множители и $g_\lambda=a_\lambda$ для каждого собственного значения $\lambda$. Так как корни характеристического многочлена --- это ровно собственные значения, сумма всех алгебраических кратностей равна степени многочлена: $n=\sum_{\lambda} a_\lambda$. По условию $a_\lambda=g_\lambda$, значит, $n=\sum_{\lambda} g_\lambda$. По критерию диагонализируемости через геометрические кратности, это условие означает, что оператор $\alpha$ диагонализируем.
\end{proof}

\begin{Corollary}{Случай простых корней}{}
	Если характеристический многочлен $\chi_\alpha$ имеет $n$ различных корней в поле $K$,
	то оператор $\alpha$ диагонализируем.
\end{Corollary}

\begin{proof}
	В этом случае все алгебраические кратности равны $1$.
	Для каждого собственного значения $\lambda$ имеем
	\[
	1\le g_\lambda\le a_\lambda=1,
	\]
	поэтому $g_\lambda=a_\lambda=1$.
	Значит, выполнен критерий диагонализируемости через алгебраические кратности.
\end{proof}